\subsection{Further restrictions on the conjugate period-ratio gate}

\begin{result}{A.3.5}{Rational-function relations and the symmetric residual equation}{Proved}
For a non-CM elliptic curve $E/\Qb$,
\[
 \bar\tau\in\Q(\tau)\quad\Longleftrightarrow\quad E\sim\bar E.
\]
For $E\not\sim\bar E$, every algebraic dependence of $\tau,\bar\tau$ is generated by an irreducible polynomial
\[
 P(X,Y)\in\Q[X,Y],\qquad P(Y,X)=P(X,Y),\qquad
 \deg_XP=\deg_YP=d\ge2.
\]
Writing $F=\Q(\tau,\bar\tau)$, its conjugation-fixed field is
\[
 F^{\langle c\rangle}=\Q(\tau+\bar\tau,\tau\bar\tau),
 \qquad [F:F^{\langle c\rangle}]=2.
\]
\end{result}
\begin{proof}
By Schneider's theorem in A.3.1, $\tau$ is transcendental. If $\bar\tau=R(\tau)$ for $R\in\Q(X)$, conjugation gives $\tau=R(R(\tau))$. Evaluation at $\tau$ is injective on $\Q(X)$, so $R\circ R=X$. The degree of a composite rational map of projective lines is the product of the degrees; hence $\deg R=1$. Write $R(X)=(aX+b)/(cX+d)$ with rational entries and nonzero determinant. Multiplication by $c\tau+d$ puts the lattices generated by $1,\bar\tau$ and $1,\tau$ in commensurable rational spans. Clearing denominators gives an analytic isogeny, which descends to $\Qb$ by the torsion-kernel quotient argument in A.3.2. That result also proves the reverse implication.

In the dependent case, the kernel of $\Q[X,Y]\to\C$, $(X,Y)\mapsto(\tau,\bar\tau)$, is a height-one prime of a UFD, hence generated by an irreducible $P$. Conjugation preserves this kernel and exchanges the variables, so $P(Y,X)=aP(X,Y)$ for a nonzero rational $a$. Applying this twice gives $a=\pm1$. The negative case makes $P$ divisible by $X-Y$; irreducibility would then force $\tau=\bar\tau$, a contradiction. The polynomial is therefore symmetric, and its separate degrees agree. Degree zero contradicts transcendence; degree one expresses one coordinate as a rational function of the other and is excluded by the first assertion.

Finally put $r=\tau+\bar\tau$, $s=\tau\bar\tau$. Both are real and fixed by conjugation, and $\tau$ satisfies $T^2-rT+s=0$, while $\bar\tau=r-\tau$. Thus $[F:\Q(r,s)]\le2$. It is exactly two because $\tau$ is not real. The unique nontrivial automorphism is conjugation, proving the fixed-field assertion.
\end{proof}

The rational coefficient field matters. For $R\in\Qb(X)$ the same argument gives $\bar R\circ R=X$ and degree one, but the resulting matrix has algebraic entries. The lattice argument above specifically uses rational entries; no unproved passage between these coefficient fields is used.

\begin{result}{A.3.6}{The exact product-period quotient and tensor relation}{Proved reductions; arithmetic gate retained}
Let $E/\Qb$ be nonisogenous to $\bar E$, and put $A=E\times\bar E$. In the absolute-Hodge abelian category, its period torsor has group
\[
 \Delta_2=\{(g_1,g_2)\in\GL_2^2:\det g_1=\det g_2\},\qquad \dim\Delta_2=7.
\]
It has a surjective algebraic flag map
\[
 q:P_A\longrightarrow\mathbf P^1\times\mathbf P^1,
 \qquad p\longmapsto(p^{-1}F^1_{E,\mathrm{dR}},p^{-1}F^1_{\bar E,\mathrm{dR}}),
\]
with fibres of dimension five and
\[
 q(p_{\mathrm{an}})=([1:\tau],[1:-\bar\tau]).
\]
Thus the named instance of Bost--Charles's period conjecture for $h^1(E_8)\oplus h^1(\bar E_8)$ implies the unresolved degree-two assertion for A.3.4 through this map. The exact residual assertion is injectivity of
\[
 \Qb[X,Y]\longrightarrow\C,\qquad X\longmapsto\tau,\quad Y\longmapsto-\bar\tau.
\]
A proposed polynomial relation $P=\sum c_{ij}X^iY^j$ of bidegree $(m,n)$ is equivalently the period relation
\[
 \sum_{i,j}c_{ij}\omega_1^{m-i}\omega_2^i
 \bar\omega_1^{\,n-j}\bar\omega_2^{\,j}=0
 \tag{A.3.6}
\]
in $\operatorname{Sym}^m H^1(E)\otimes\operatorname{Sym}^n H^1(\bar E)$, with the unaltered conjugate coordinate in this last formula.
\end{result}
\begin{proof}
The Hodge group projects onto each $\SL_2$. Its connected derived group is a subdirect product of these two simple groups. The Lie-algebra Goursat argument gives either the full product or the graph of an isomorphism. In the graph case, an automorphism of $\mathfrak{sl}_2$ is induced by $\mathrm{PGL}_2(\Q)$, giving an intertwiner of the two standard rational representations. The Hodge group has no positive-dimensional centre, since its central image in each $\SL_2$ is finite. The intertwiner consequently commutes with the Hodge circle, giving a rational Hodge isomorphism of the two $H^1$'s. The lattice construction of A.3.2 then gives an isogeny, contrary to the hypothesis. Hence the Hodge group is $\SL_2^2$. Adding the common weight character gives $\Delta_2$; absolute Hodge cycles on abelian varieties identify this with the motivic group, as used in A.4.9.

The map $q$ is algebraic because both de Rham lines are algebraic. Its equivariance and the transitivity of $\SL_2^2$ on the target prove surjectivity and the fibre dimension. On the conjugate curve use the oriented period basis $(\bar\omega_1,-\bar\omega_2)$, which proves the displayed value. A generic point of a torsor maps to a generic point under a dominant morphism: the nonzero pullback of any nonzero function on a target chart cannot vanish there. This derives the stated implication from \cite[Conjecture 2.12]{GenBostCharles} for the named motive. Independence over $\Qb$ and over $\Q$ agree, by multiplying a relation over a finite number field by its conjugates.

Multiplication by $\omega_1^m\bar\omega_1^n$ proves (A.3.6). K\"unneth identifies its monomials with integrals of the product holomorphic form over product cycles on $E^m\times\bar E^n$. The ordinary symmetric projector on the tensor product of $H^1$ Hodge structures gives the specified symmetric powers; in the geometric cohomology projector the permutation signs are included. One may equivalently retain the full K\"unneth tensor product. Its symmetric summand has Hodge types $(m+n,0)$ and $(0,m+n)$ and thus level $m+n$. Tate twisting preserves this level. A pure constituent of a 1-motive has level at most one, so for the residual $m=n=d\ge2$ there is no isomorphism, even after Tate twist, to such a constituent. This proves the precise obstruction to applying the linear 1-period theorem directly to (A.3.6); it makes no claim about every possible extension of its methods.
\end{proof}

\paragraph{A compatibility test retaining the individual maximal degrees.}
Choose real $t>0,r,s,p,h$ algebraically independent over $\Qb(\pi)$ and put
\[
 z=t+it^2,\quad u=r+is,\quad v=p+ih,\quad
 (\omega_1,\omega_2,\eta_1,\eta_2)
 =(u,uz,v,zv-2\pi i/u).
\]
Take actual complex conjugates for the second tuple. Legendre's relation holds, each four-entry field has transcendence degree four, and the combined field has degree six: it recovers exactly $\Qb(\pi,r,s,t,p,h)$. The ten entries consisting of $1,2\pi i$ and the two four-entry tuples are $\Qb$-linearly independent. Indeed comparison of the independent $p,h$ coefficients first gives $a+bz=0$ and $a'+b'\bar z=0$, forcing those four algebraic coefficients to vanish; comparison of $r,s$ then does the same for the holomorphic entries, and transcendence of $\pi$ finishes. Nevertheless
\[
 4(z-\bar z)^2+(z+\bar z)^4=0,
 \qquad\trdeg_{\Q}\Q(z,\bar z)=1.
\]
These are explicitly constructed complex numbers; no algebraic invariants or algebraic $j(z)$ have been asserted. The construction proves that the individual maximal period degrees, linear product-period independence, the determinant relations and conjugation leave an additional arithmetic step. The symbolic identities and their inputs are in \texttt{check\_tau\_gate.py}.

\paragraph{Audit of further theorem applications.}
Huber--W\"ustholz's \cite[Theorems 9.7 and 9.10, pp.78--79]{TauHuberWuestholz} controls 1-periods; the exact tensor and Hodge-level map above records the missing input for this use. Pila--Tsimerman's \cite[Theorems 1.1--1.3, pp.1--3]{TauPilaTsimerman} applied to the graph of $j^2$ at $(\tau,-\bar\tau,j_0,\bar j_0)$ has a constant moduli projection, which is a weakly special point. The nonconstant-projection hypothesis therefore fails. Its differential version imposes $D(j(\tau))=j'(\tau)D(\tau)$; algebraicity of $j_0\ne0,1728$ makes a compatible $\Q$-derivation kill $\tau$ as well. An arbitrary field derivation cannot be assumed to satisfy that analytic chain rule.

The functional period theorem of \cite[Theorems 1.1 and 1.3]{EllBakkerTsimerman2025} uses the relative group over a complex base; the constant product has trivial relative group. The germ-field identification of \cite[Theorem 1.1 and \S4.2]{TauBakkerPilaTsimerman} does not give injective specialization at a point: a local parameter $s-s_0$ is a nonzero germ with zero value, and evaluation does not commute with its derivative. The complete source audit, the G-function matrix $G(s)=P(s)P(s_0)^{-1}$ and the additional global-relation requirements are recorded in \texttt{agent\_tau\_gate.md}. No result checked there excludes the remaining symmetric equations of bidegree $(d,d)$, $d\ge2$, for the specified $E_8$.
